\chapter{Evaluation norms and operator closure}

A residue value is a linear functional on compact smooth tests.
Completing those tests in the squared sum of their residue values
changes the Hilbert space. The change can be measured exactly for
finitely many zeros, without any estimate on the global zero set.

\section{Finite evaluation and its graph}

\begin{lemma}[independence of exponential polynomials]
\label{mo:exponential-polynomials}
Let $\lambda_1,\ldots,\lambda_n$ be distinct complex numbers.
If polynomials $p_j$ satisfy
$\sum_j p_j(t)e^{\lambda_jt}=0$ on an interval, then all $p_j$ vanish.
A nonzero sum of this form does not belong to $L^2(\mathbb R,dt)$.
\end{lemma}
\begin{proof}
For independence, fix $j$ and apply
$\prod_{k\ne j}(\partial_t-\lambda_k)^{d_k+1}$, where
$d_k\geq\deg p_k$. All terms except the $j$th vanish. On its
polynomial factor this operator is
$\prod_{k\ne j}(\partial_t+\lambda_j-\lambda_k)^{d_k+1}$.
Each factor is invertible on a space of polynomials of bounded
degree: in the monomial basis it is triangular with nonzero
constant diagonal. Therefore $p_j=0$.

For the $L^2$ assertion, choose the largest real part $a$ among
nonzero terms, and the largest polynomial degree $d$ among terms
with that real part. On intervals $[T,T+L]$, divide the sum by
$e^{aT}T^d$. As $T\to+\infty$, its leading part consists of a fixed
finite family $e^{(a+ib_j)u}$, $0\leq u\leq L$, with coefficients
of fixed nonzero magnitudes and phases $e^{ib_jT}$. The Gram
matrix of this family on $[0,L]$ is positive definite by the first
part. All lower-degree or lower-real-part terms tend to zero
uniformly after division. Thus, for sufficiently large $T$,
\[
 \int_T^{T+L}\left|\sum_jp_j(t)e^{\lambda_jt}\right|^2dt
 \geq C e^{2aT}T^{2d}
\]
for some $C>0$. If $a\geq0$, summing over disjoint such intervals
proves divergence. If $a<0$, replace $t$ by $-t$; the largest
real part then becomes positive. This proves the claim.
\end{proof}

\begin{theorem}[finite residue graph]
\label{mo:finite-graph}
For every real $c$ and every nonempty finite zero set $Z$, the map
$\mathcal R_Z:\mathscr D\to\mathscr J_Z$ is onto. Its graph is
dense in $\mathscr H_c\oplus\mathscr J_Z$, where $\mathscr J_Z$
has its Euclidean norm. Hence it is unbounded and not closable
as an operator from $\mathscr H_c$. Its kernel is dense in
$\mathscr H_c$.
\end{theorem}
\begin{proof}
Write $g=U_cf$. The Taylor functionals before the invertible
transformation \eqref{mo:jet-formula} are
\[
 F^{(k)}(\rho)=\int_{\mathbb R}t^k
                       e^{(\rho-c)t}g(t)\,dt.
\]
The kernels are linearly independent by
Lemma~\ref{mo:exponential-polynomials}. No nonzero functional on the
finite-dimensional target can annihilate the range. Thus the
range is the whole target.

Suppose $(u,b)\in L^2(\mathbb R)\oplus\mathscr J_Z$ is orthogonal
to the graph in logarithmic coordinates. For every compact smooth
$g$ this says
\[
 \int g(t)\overline{u(t)}\,dt+\int g(t)q_b(t)\,dt=0,
\]
where $q_b$ is a fixed linear combination of the exponential
polynomials above. Thus $q_b=-\overline u$ as a distribution,
so $q_b\in L^2$. The lemma forces $q_b=0$. Independence and the
invertible jet transformation give $b=0$, and then $u=0$.
A subspace with zero orthogonal complement is dense. The graph
closure is therefore the entire product, which contains
$\{0\}\oplus\mathscr J_Z$ and cannot be a graph.

Choose test functions lifting a basis of $\mathscr J_Z$; their
linear extension is a right inverse $G:\mathscr J_Z\to\mathscr D$.
As a map into $\mathscr H_c$, $G$ is bounded because its domain is
finite-dimensional. Graph density gives, for each $f\in\mathscr H_c$,
$g_n\to f$ and $\mathcal R_Zg_n\to0$. Then
$g_n-G\mathcal R_Zg_n$ lies in the kernel and converges to $f$.
\end{proof}

Define a positive semidefinite form on $\mathscr D$ by
\[
 q_Z(f,h)=\sum_{\rho\in Z}\sum_{k=0}^{r_\rho-1}
 R_{\rho,k}\mathcal Mf\;
 \overline{R_{\rho,k}\mathcal Mh}.
\]
Its nullspace is exactly $\ker\mathcal R_Z$. On the quotient
$\mathscr D/\ker\mathcal R_Z$, this is an inner product.
Surjectivity in Theorem~\ref{mo:finite-graph} identifies its completion
isometrically with $\mathscr J_Z$. The nullspace has infinite
dimension, since $\mathscr D$ has infinite dimension and the
quotient has finite dimension. Thus this positive norm does not
give an injective realization of the original function space.

The form is not closable in $\mathscr H_c$. For a positive form,
closability requires $q(f_n)\to0$ whenever $f_n\to0$ in the
original Hilbert norm and $q(f_n-f_m)\to0$. Dense graph gives
$f_n\to0$ with $\mathcal R_Zf_n$ tending to a fixed nonzero
vector. This sequence violates the form criterion. The same
argument applies to the countable form below whenever its stated
summability hypothesis holds.

\begin{theorem}[symmetry in the finite evaluation norm]
\label{mo:finite-symmetry}
The dilation operator descends to the quotient with form $q_Z$ and
extends to the matrix $K_c$ of \eqref{mo:jet-generator}.
It is symmetric for the Euclidean jet norm exactly when all zeros
in $Z$ are simple and satisfy $\operatorname{Re}\rho=c$.
If either condition fails, no positive definite inner product on
the full jet target can make $K_c$ symmetric.
\end{theorem}
\begin{proof}
The identity $\mathcal R_ZP_c=K_c\mathcal R_Z$ makes the kernel
invariant and proves the induced action. In finite dimension the
induced operator is everywhere defined, bounded, and closed.
If every block has size one and real eigenvalue, its matrix is
real diagonal and hence symmetric.

Conversely, symmetry on a positive inner-product space forces
an eigenvalue $z$ to be real, since
$\langle Kv,v\rangle=z\|v\|^2$ must be real for an eigenvector.
Thus $\operatorname{Re}\rho=c$. A nontrivial Jordan block then
has vectors $u,v\ne0$ with $(K-z)u=v$ and $(K-z)v=0$.
Symmetry would give
$\|v\|^2=\langle(K-z)u,v\rangle=
\langle u,(K-z)v\rangle=0$, a contradiction.
\end{proof}

For one double zero on the line, the matrix is
\[
 \begin{pmatrix}\gamma&-i\\0&\gamma\end{pmatrix}.
\]
Its two real eigenvalues do not remove its nilpotent part. A
positive scalar-residue norm using only $R_{\rho,0}$ generally
has a different defect: its kernel need not be invariant under
$P_c$, because the omitted coordinate $R_{\rho,1}$ appears in
its derivative. The full jets are required before asking for
an induced dilation operator.

\section{Countable simple zeros and weighted evaluations}

Let $\rho_j=c+i\gamma_j$ be a countable set of distinct simple
zeros, with no finite accumulation. Choose positive weights $w_j$.
The raw residue norm would be
\begin{equation}
 q_w(f)=\sum_j w_j
       \left|\frac{F(\rho_j)}{\Xi'(\rho_j)}\right|^2.
 \label{mo:weighted-residue-norm}
\end{equation}
Put $\mu_j=w_j/|\Xi'(\rho_j)|^2$. Multiplication of the residue
coordinates by $\Xi'(\rho_j)$ identifies the target norm with
$\ell^2(\mu)$, whose norm is $\sum_j\mu_j|a_j|^2$.
This conversion retains the residue normalization, including its
complex phase.

A sufficient summability condition is
\begin{equation}
 \sum_j\mu_j(1+\gamma_j^2)^{-b}<\infty
 \quad\text{for some }b\geq0.
 \label{mo:weight-condition}
\end{equation}
It is a hypothesis on the chosen weights and the actual divisor.
It is not inferred for the raw choice $w_j=1$. Under
\eqref{mo:weight-condition}, define
\[
 E:\mathscr D\longrightarrow\ell^2(\mu),\qquad
 (Ef)_j=F(c+i\gamma_j),\qquad
 (Ka)_j=\gamma_ja_j.
\]
The maximal domain of $K$ consists of sequences with
$\sum_j\mu_j\gamma_j^2|a_j|^2<\infty$. Compact smooth tests have
Schwartz Mellin transforms on the vertical line, by repeated
integration by parts in $t$. Therefore $Ef\in\operatorname{Dom}K$
and $EP_c=KE$ on $\mathscr D$.

\begin{theorem}[the atomic completion and the original norm]
\label{mo:countable-completion}
Under \eqref{mo:weight-condition}, $E\mathscr D$ is dense in
$\ell^2(\mu)$ and is a graph core for $K$.
The completion of $\mathscr D/\ker E$ in the norm $\|Ef\|$
is therefore $\ell^2(\mu)$. The induced dilation has self-adjoint
closure $K$ on its stated maximal domain.

As a map from the original space $\mathscr H_c$, $E$ has dense
graph in $\mathscr H_c\oplus\ell^2(\mu)$. In particular it is
not closable and has no closable extension agreeing on
$\mathscr D$.
\end{theorem}
\begin{proof}
Fix $j_0$ and use $f_L$ from \eqref{mo:continuous-sequence} with
$\gamma_0=\gamma_{j_0}$. Put
$\Psi(y)=\int\psi(t)e^{iyt}\,dt$. Then
\[
 (Ef_L)_j=\Psi\bigl(L(\gamma_j-\gamma_{j_0})\bigr).
\]
The central value is one and every other value tends to zero.
For $L\geq1$, rapid decay of $\Psi$ gives a uniform bound
$C_n(1+|\gamma_j-\gamma_{j_0}|)^{-n}$ for every $n$.
Choosing $n$ sufficiently large in
\eqref{mo:weight-condition} gives a summable majorant even after
multiplication by $1+\gamma_j^2$. Dominated convergence proves
\[
 Ef_L\longrightarrow e_{j_0},\qquad
 KEf_L\longrightarrow\gamma_{j_0}e_{j_0}
 \quad\text{in }\ell^2(\mu),
\]
while $f_L\to0$ in $\mathscr H_c$.
Finite-coordinate vectors form a graph core for the maximal
real diagonal operator $K$, by truncation of the two convergent
weighted sums. The approximations just constructed show that
$E\mathscr D$ is also a graph core. The diagonal adjoint-domain
calculation proves self-adjointness, as in
Lemma~\ref{mo:mellin-unitary}.

The graph closure of $E$ contains $(0,e_j)$ for every $j$, hence
$\{0\}\oplus\ell^2(\mu)$. It also contains $(f,Ef)$ for every
test $f$. Subtracting the vertical vectors and using density of
$\mathscr D$ in $\mathscr H_c$ gives the whole product. This
proves the last assertion.
\end{proof}

The hypotheses of this theorem can always be imposed on a selected
countable critical-line simple divisor by changing the weights.
For example, fix an enumeration and set
\[
 \mu_j=2^{-j}(1+\gamma_j^2)^{-j},\qquad
 w_j=|\Xi'(\rho_j)|^2\mu_j.
\]
All polynomial moments of $\mu$ are finite: for a fixed degree,
only finitely many initial terms need separate consideration.
This produces a positive evaluation completion. The weights
are chosen data; they are not the raw residue weights or a native
arithmetic pairing.

For $w_j=1$, the exact domain of \eqref{mo:weighted-residue-norm}
is
\[
 \mathscr D_{\mathrm{raw}}=
 \left\{f\in\mathscr D:
 \sum_j|F(\rho_j)/\Xi'(\rho_j)|^2<\infty\right\}.
\]
No assertion that this equals $\mathscr D$, is dense in
$\mathscr H_c$, or satisfies \eqref{mo:weight-condition} is made
without bounds on the inverse derivatives. Even existence of an
individual scalar residue does not imply convergence of this sum.
If \eqref{mo:weight-condition} is established for the raw weights,
Theorem~\ref{mo:countable-completion} applies and still gives
nonclosability in the original norm. Without that estimate, a
full vector-valued closure claim remains separate from the
finite-coordinate obstruction.

\begin{proposition}[bounded spectral transfer]
\label{mo:bounded-transfer}
A bounded map intertwining the unitary groups of $P_c$ on
$\mathscr H_c$ and $K$ on $\ell^2(\mu)$ is zero, in either
direction.
\end{proposition}
\begin{proof}
In the direction from the atomic space, the image of $e_j$ would
be an eigenvector of $P_c$ with eigenvalue $\gamma_j$.
Under $\mathcal U_c$, group covariance for all rational times
forces its support to lie in $\{\gamma_j\}$, a null set.
Thus it is zero, and density of the finite-coordinate vectors
finishes the argument. Taking the adjoint gives the opposite
direction, since both groups are unitary.
\end{proof}

\begin{proposition}[occupation symmetry in a sampling norm]
\label{mo:occupation-positive-obstruction}
Let $\rho_j=c+i\gamma_j$ be a nonempty countable set of simple
zeros, with positive weights $\mu_j$ and
$\sum_j\mu_j<\infty$. The evaluation form
\[
 q(a,b)=\sum_j\mu_jF_a(\rho_j)\overline{F_b(\rho_j)}
 \quad(a,b\in\mathscr E)
\]
is well-defined and positive semidefinite. The occupation operator
$N$ is not symmetric for this form on its finite-support core.
\end{proposition}
\begin{proof}
Every finite Dirichlet polynomial is bounded on a fixed vertical
line. The summability assumption therefore defines the form.
Put $\phi(t)=\sum_j\mu_je^{i\gamma_jt}$. Uniform convergence
makes $\phi$ continuous, and $\phi(0)=\sum_j\mu_j>0$.
The coordinate pairing is
\[
 q(v_m,v_n)=(mn)^{-c}\phi(\log(n/m)).
\]
Symmetry would require
$(\log m-\log n)q(v_m,v_n)=0$ for all $m,n$.
Taking $n=m+1$ would give
$\phi(\log(1+1/m))=0$. Continuity at zero contradicts
$\phi(0)>0$.
\end{proof}

Thus even a positive critical-line evaluation completion cannot
make occupation energy symmetric on this core. The same norm
does make the induced continuous dilation self-adjoint under
Theorem~\ref{mo:countable-completion}. These statements concern different
source operators, with the exact actions in \eqref{mo:two-actions}.

\section{A countable completion with all principal parts}

The full principal parts also admit a Hilbert completion when
multiplicities are unbounded or zeros lie off the chosen line.
The construction below chooses the norm after the divisor and
finite interpolation functions have been specified.

Enumerate distinct zeros as $\rho_j$, $j\geq1$, with
$|\rho_j|\to\infty$, and let $r_j$ be their finite multiplicities.
Fix real $c$. Write $R_jf\in\mathbb C^{r_j}$ for the full residue
vector of $\mathcal Mf$ and set
\[
 z_j=i(c-\rho_j),\qquad K_j=z_jI-iS_j,
 \qquad (S_ja)_k=a_{k+1},\quad a_{r_j}=0.
\]
Each block has its Euclidean norm, so $\|S_j\|\leq1$.
For positive numbers $w_j$ define
\[
 \mathscr J_w=\left\{a=(a_j):
              \sum_jw_j\|a_j\|^2<\infty\right\},\qquad
 (K_wa)_j=K_ja_j.
\]
Its maximal domain is
\begin{equation}
 \operatorname{Dom}K_w=
 \left\{a\in\mathscr J_w:
             \sum_jw_j|z_j|^2\|a_j\|^2<\infty\right\}.
 \label{mo:full-domain}
\end{equation}
Indeed $K_w$ differs from multiplication by $z_j$ by the bounded
operator $-i\bigoplus_jS_j$. Equivalently its domain requires
$\sum_jw_j\|K_ja_j\|^2<\infty$.

\begin{theorem}[weighted full-jet completion]
\label{mo:all-jets}
For every such zero divisor and every real $c$, positive weights
$w_j$ can be chosen with the following properties. The map
$R:f\mapsto(R_jf)_j$ sends $\mathscr D$ into
$\operatorname{Dom}K_w$ and satisfies $RP_c=K_wR$.
Its range is dense and is a graph core for $K_w$.
The completion of $\mathscr D/\ker R$ in the norm $\|Rf\|$
is $\mathscr J_w$, and the induced dilation closes to $K_w$.
The graph of $R$ in $\mathscr H_c\oplus\mathscr J_w$ is dense.
Thus $R$ and its positive evaluation form are not closable in
the original norm.

The closed operator $K_w$ has compact resolvent and spectrum
$\{z_j:j\geq1\}$. The generalized eigenspace at $z_j$ is one
block of dimension $r_j$. It is self-adjoint exactly when all
$r_j=1$ and $\operatorname{Re}\rho_j=c$.
\end{theorem}
\begin{proof}
Let $I_n=\{(j,k):j\leq n,\ 0\leq k<r_j\}$ and denote its
coordinate basis by $e_h$. The finite graph theorem supplies,
for each $n$ and $h\in I_n$, a test $f_{n,h}$ such that
\begin{equation}
 \|f_{n,h}\|_{\mathscr H_c}\leq n^{-1},\qquad
 (R_jf_{n,h})_{j\leq n}=e_h.
 \label{mo:small-lifts}
\end{equation}
Here exact interpolation follows from graph density by correction
with the bounded finite-dimensional right inverse used in its proof.
All these choices precede the choice of weights.

For $g=U_cf$ supported in $[-m,m]$, each residue is an integral
against an exponential polynomial. Hence there is a finite
constant $C_{j,m}$ with
$\|R_jf\|\leq C_{j,m}\|g\|_2$ on that support.
One explicit choice is
\[
 C_{j,m}^2=\sum_{k=0}^{r_j-1}\int_{-m}^m
 \left|e^{(\rho_j-c)t}
 \sum_{\ell=0}^{r_j-k-1}
 b_{\rho_j,-k-\ell-1}\frac{t^\ell}{\ell!}\right|^2dt.
\]
Cauchy--Schwarz proves the bound. Define the finite maximum
\[
 A_j=\max\left(
 1,\ \max_{1\leq m\leq j}C_{j,m}^2,
 \max_{\substack{1\leq n<j\\h\in I_n}}
                  n^2\|R_jf_{n,h}\|^2\right),
 \qquad
 w_j=\frac{2^{-j}}{(1+|z_j|^2)A_j}.
\]
The last maximum is omitted for $j=1$.
For any fixed compact smooth test and sufficiently large $j$,
the summand $w_j(1+|z_j|^2)\|R_jf\|^2$ is at most
$2^{-j}\|U_cf\|_2^2$. This proves membership in the domain.
The coordinate identity from Proposition~\ref{mo:jet-comparison} therefore
holds as an identity in $\mathscr J_w$.

For a fixed coordinate $h$, take $n$ large enough that $h\in I_n$.
The first $n$ blocks of $Rf_{n,h}-e_h$ vanish. The remaining
blocks satisfy
\[
 \sum_{j>n}w_j(1+|z_j|^2)\|R_jf_{n,h}\|^2
 \leq n^{-2}\sum_{j>n}2^{-j}.
\]
Since $\|K_j\|\leq |z_j|+1$, this gives
$Rf_{n,h}\to e_h$ and $K_wRf_{n,h}\to K_we_h$.
Finite block truncations are a graph core for the maximal operator,
by convergence of the two sums defining its graph norm.
Thus $R\mathscr D$ is a graph core and is dense.
Equation \eqref{mo:small-lifts} also gives $f_{n,h}\to0$.
The graph closure contains $\{0\}\oplus\mathscr J_w$.
Subtracting its vertical vectors from $(f,Rf)$ and using density
of $\mathscr D$ proves density in the whole product.
The quotient and nonclosability assertions follow.

The maximal block operator is closed: convergence in its graph
implies convergence in every finite block, where each matrix is
closed. Its adjoint is the maximal operator with blocks
$\overline z_jI+iS_j^*$, as testing against finite block vectors
shows. These formulas also prove self-adjointness when all blocks
are real scalars. If it is symmetric, its restriction to every
finite block is symmetric. Theorem~\ref{mo:finite-symmetry} forces every
block to be a real scalar.

Choose $\lambda$ different from every $z_j$. For all sufficiently
large $j$, $|z_j-\lambda|>2$. The finite geometric series for
$(K_j-\lambda)^{-1}$ then gives
\[
 \|(K_j-\lambda)^{-1}\|
 \leq\frac1{|z_j-\lambda|-1}\longrightarrow0.
\]
The finitely many other inverses are bounded. Their direct sum
maps $\mathscr J_w$ into \eqref{mo:full-domain}, since
$K_ja_j=y_j+\lambda a_j$ for
$a_j=(K_j-\lambda)^{-1}y_j$.
Thus $\lambda$ belongs to the resolvent. Truncating this inverse
to finitely many finite-dimensional blocks converges in operator
norm, so the resolvent is compact. Each $z_j$ is an eigenvalue.
On every other block $K_l-z_j$ is invertible, so its generalized
eigenspace is precisely the $j$th block.
\end{proof}

The theorem requires no bound on the multiplicities. The norm
depends on the Laurent coefficients and the chosen interpolation
functions. Making a weight smaller preserves every estimate and
every graph approximation in the proof. For a divisor invariant
under $\rho\mapsto2c-\overline\rho$, one may therefore replace
the two weights on each orbit by their minimum. The resulting
residue involution is antiunitary. This does not make a nontrivial
Jordan block symmetric.

For a double zero on the line, the ordinary trace of a polynomial
in $K_j$ is $2p(\gamma_j)$, while the operator retains its
nonzero nilpotent part. More generally,
\[
 \operatorname{tr}p(K_j)=r_jp(z_j).
\]
This follows by triangularity. Consequently these traces cannot
distinguish the full block from $r_j$ repeated scalar eigenvalues.
Equality of spectral traces alone does not construct a map of
their full jet representations.


\section{Stronger source norms and boundary values}

Evaluation can be continuous after the source topology changes.
On finite occupation sequences define, for real $a$,
\[
 \|u\|_a^2=\sum_{m\geq1}(1+(\log m)^2)^a|u_m|^2,
\]
and let $\mathscr H_{\mathrm{occ},a}$ be the corresponding weighted
$\ell^2$ space. Fix $\rho=1/2+i\gamma$ and an integer $k\geq0$.

\begin{proposition}[the exact logarithmic weight threshold]
\label{mo:weight-threshold}
The functional
$u\mapsto F_u^{(k)}(\rho)=\sum_m u_m(-\log m)^k m^{-\rho}$
on finite sequences is bounded for $\|\cdot\|_a$ exactly when
$a>k+1/2$. If $\operatorname{Re}\rho>1/2$, it is bounded for
every real $a$. If $\operatorname{Re}\rho<1/2$, no real $a$
makes it bounded.
\end{proposition}
\begin{proof}
The squared dual norm of a finite-coordinate functional increases
under truncation to
\[
 \sum_m\frac{(\log m)^{2k}m^{-2\operatorname{Re}\rho}}
                    {(1+(\log m)^2)^a}.
\]
Cauchy--Schwarz proves sufficiency. Truncated coefficient vectors
proportional to the weighted conjugate kernel attain the finite
dual norms, proving necessity. At real part $1/2$, integral
comparison with $t=\log x$ gives the tail integral
$\int^\infty t^{2k-2a}\,dt$. It is finite exactly when
$a>k+1/2$. To the right, exponential decay in $t$ dominates every
power. To the left, exponential growth defeats every such power.
\end{proof}

At a zero of order $r$, the scalar residue functional on finite
occupation sequences has highest derivative order $r-1$, with
nonzero coefficient. On the critical line its exact threshold
is therefore $a>r-1/2$. The same condition controls all full jet
coordinates. Lower derivative terms cannot cancel the nonzero
leading polynomial in $\log m$.

These extensions are boundary-value functionals of the Dirichlet
series from Proposition~\ref{mo:dirichlet-domain}. Their defining sums converge
absolutely under the stated weight condition. They do not supply
holomorphic continuation across the boundary. Calling them actual
meromorphic residues for an infinite sequence additionally requires
such continuation. Nor do stronger norms repair
Proposition~\ref{mo:occupation-defect}, which already fails on finite sequences.

The occupation operator remains self-adjoint in
$\mathscr H_{\mathrm{occ},a}$ on its maximal weighted domain,
because both the norm and the operator are diagonal in $v_m$.
For continuous tests, a different change is needed to bound point
evaluation: put $\|g\|_{a,t}^2=\int(1+t^2)^a|g(t)|^2dt$.
The same dual-integral argument bounds the $k$th Fourier evaluation
exactly when $a>k+1/2$. This norm is not translation-invariant.
For $w(t)=(1+t^2)^a$ and compact smooth $g,h$,
\[
 \langle i g',h\rangle_w-\langle g,i h'\rangle_w
 =-i\int w'(t)g(t)\overline{h(t)}\,dt.
\]
Thus the same dilation generator need not remain symmetric in
this stronger continuous norm. A topology that controls residues
must be checked against the required operator and pairing.
